\documentclass[11pt,a4paper]{article}
\usepackage[T1]{fontenc}
\usepackage[utf8]{inputenc}
\usepackage[margin=2.2cm]{geometry}
\usepackage{lmodern}
\usepackage{amsmath}
\usepackage{graphicx}
\usepackage{booktabs}
\usepackage{textcomp}
\usepackage[hidelinks]{hyperref}
\usepackage[numbers,sort&compress]{natbib}
\graphicspath{{../results/figures/}}

\title{Where Was This Photo Taken?\\[2pt]
\large Classical vs.\ Learned Feature Matching for Visual Localization under Image Degradations}
\author{Ozan Taskin\\\small Sorbonne Université, M1 Computer Science (IMA)}
\date{November 2026}

\begin{document}
\maketitle

\begin{abstract}
%ABSTRACT%
\end{abstract}

\section{Introduction}
Visual localization estimates the 6-DoF pose of a camera from a single image, given a
3D map of the scene. It is a core building block of robot navigation, augmented reality
and autonomous driving. Structure-based pipelines establish 2D--3D correspondences between
the query image and the map, and then solve a Perspective-n-Point (PnP) problem inside
RANSAC~\cite{sarlin2019coarse}. The quality of those correspondences, and hence of the pose,
depends on the local features and the matcher.

Learned features and matchers, such as SuperPoint~\cite{detone2018superpoint},
ALIKED~\cite{zhao2023aliked} and LightGlue~\cite{lindenberger2023lightglue}, are reported to
outperform handcrafted SIFT~\cite{lowe2004sift} on challenging benchmarks. This project asks a
narrower, practical question: \emph{on a small outdoor scene, how much do learned matchers gain
over SIFT, and does the gap widen when query images are degraded (low light, blur, low
resolution)?} The answer matters for robots and phones whose cameras rarely see the scene in
the conditions under which the map was built.

\paragraph{Contributions.}
(i) A compact, reproducible benchmark built on hloc~\cite{sarlin2019coarse} and
COLMAP~\cite{schoenberger2016sfm}; (ii) a comparison of three feature/matcher pairs under four
controlled query conditions with identical maps, retrieval and solver; (iii) an analysis of
when and why localization fails.

\section{Method}
\paragraph{Map building.}
For each method we extract local features on all database images and match each image to
its 10 most co-visible neighbours in the reference model. Matches are verified with the
known relative poses (epipolar check) and triangulated with the ground-truth camera poses,
as in hloc's Cambridge pipeline. Each method therefore gets its own 3D map, built from its
own features, while camera poses stay fixed.

\paragraph{Localization.}
For each query, NetVLAD~\cite{arandjelovic2016netvlad} retrieves the 10 most similar database
images. Query keypoints are matched to these images; the matched database keypoints carry 3D
points, giving 2D--3D correspondences. The pose is estimated with PnP inside LO-RANSAC
(12\,px threshold) followed by non-linear refinement, using the known query intrinsics.

\paragraph{Compared methods.}
\begin{itemize}\setlength\itemsep{0pt}
  \item \textbf{SIFT + NN ratio}: DoG keypoints with RootSIFT descriptors (COLMAP), mutual
        nearest neighbours with Lowe's ratio test (0.8). The classical baseline.
  \item \textbf{SuperPoint + LightGlue}: learned detector/descriptor and learned matcher with
        adaptive depth~\cite{detone2018superpoint,lindenberger2023lightglue}.
  \item \textbf{ALIKED + LightGlue}: a more recent, lightweight learned feature with the same
        matcher~\cite{zhao2023aliked}.
\end{itemize}
All methods use at most 2048 keypoints per image and images resized to 1024\,px.

\paragraph{Query degradations.}
Only query images are modified; database images, maps and parameters are unchanged
(Fig.~\ref{fig:conditions}).
\emph{Dark}: gamma 1.8, exposure gain 0.35 and Gaussian sensor noise ($\sigma=4$ grey levels).
\emph{Blur}: Gaussian blur with $\sigma=3$\,px.
\emph{Low-res}: downsampling by 4 and upsampling back.
Degradations are deterministic (fixed per-image seeds).

\begin{figure}[t]
  \centering
  \includegraphics[width=\linewidth]{conditions.png}
  \caption{One query image under the four conditions.}
  \label{fig:conditions}
\end{figure}

\section{Experimental Setup}
\paragraph{Data.} We use the ShopFacade scene of Cambridge Landmarks~\cite{kendall2015posenet}:
231 database and 103 query images of a shop front, with reference poses from a COLMAP
reconstruction distributed with hloc.

\paragraph{Metrics.} Position error is the distance between estimated and ground-truth camera
centres; rotation error is the angle of the relative rotation. We report median errors and the
percentage of queries within (25\,cm, 2\textdegree), (50\,cm, 5\textdegree) and
(5\,m, 10\textdegree). Queries where RANSAC finds no pose count as failures (infinite error).
Runtime is the wall-clock time per query for feature extraction, matching and pose estimation,
on a laptop CPU (Apple M2, 8\,GB RAM, no GPU).

%RESULTS%

\bibliographystyle{plainnat}
\bibliography{references}
\end{document}
